\SourceLeaf{282}%
is to be a minimum. The rules of the calculus of variations give the equation
\[
\frac{\mathrm{d}n}{\mathrm{d}u}\mathord{.}\BareRoot{(1+\frac{\mathrm{d}u^2}{\mathrm{d}x^2})}\mathrm{d}x=\mathrm{d}\frac{\frac{n\mathrm{d}u}{\mathrm{d}x}}{\BareRoot{(1+\frac{\mathrm{d}u^2}{\mathrm{d}x^2})}}
\]
or
\[
\frac{\mathrm{d}n}{\mathrm{d}u}\mathord{.}\frac{\mathrm{d}x}{\cos\psi}=\mathrm{d}\mathord{.}n\sin\psi
\]
Here \(\tfrac{\mathrm{d}n}{\mathrm{d}u}\) denotes the partial derivative. This formula is exact and holds in general. For our purpose, however, only the portion between \(F\) and \(G\) of the curve \(M\) comes into consideration. At all its points, \(u\) and \(\psi\) can have only very small values. We may therefore write \(1\) in place of \(\cos\psi\) and \(\operatorname{tang}\psi\) in place of \(\sin\psi\), and hence set
\[
\frac{\mathrm{d}n}{\mathrm{d}u}\mathord{.}\mathrm{d}x=\mathrm{d}\mathord{.}n\operatorname{tang}\psi
\]
or
\[
n\operatorname{tang}\psi=\int\frac{\mathrm{d}n}{\mathrm{d}u}\mathrm{d}x+\text{Const.}
\]
At the same time we may use in this formula, in place of the values that \(n\) and \(\tfrac{\mathrm{d}n}{\mathrm{d}u}\) have on the line \(M\), those at the corresponding points of the line \(N\) (for \(u=0\) or \(y=0\)), which thus agree with the values of \(\tfrac{1}{m}\) and \(-\tfrac{\mathrm{d}m}{mm\mathrm{d}u}=-\tfrac{\mathrm{d}m}{mm\mathrm{d}y}\).

To make the further developments easier, count the abscissas from the point \(F\), or set \(x=0\) at that point and at \(G\) set \(x=h\). I further set \(\tfrac{\mathrm{d}m}{m\mathrm{d}y}=l\). In general this is a function of \(x\) and \(y\), but here we consider only its value on the line \(N\) or for \(y=0\), and so as a function of \(x\) alone. Finally, let \(\psi^0,m^0,l^0\) be the particular values of \(\psi,m,l\) at \(F\), and \(\psi',m',l'\) those at \(G\). The formula above then becomes
\[
\operatorname{tang}\psi=\frac{m\operatorname{tang}\psi^0}{m^0}-m\int\frac{l}{m}\mathrm{d}x
\]
where integration starts at \(x=0\). Suppose now that \(l\) and \(m\) are expanded in the following series in successive powers of \(x\)
